% GENERATED FILE. Edit canonical lessons, then run npm run generate:books.
% canonical-source-hash: fnv1a64:18c28beb19d2d1fc
% canonical-lessons: GU-C224-daktar, GU-C224-dardi, GU-C224-nadi, GU-C224-dhabakvu, GU-C224-shvas, GU-R224-review-219-224

\chapter{At the doctor's}
\label{ch:gu-at-the-doctor-s}
\hlchaptermodality{224}

\begin{chapteropening}
\textbf{By the end of this chapter:} \emph{I can answer the doctor's questions about my pulse and breathing.}

Complete the last lesson of chapter 224: I can answer the doctor's questions about my pulse and breathing.
\end{chapteropening}

\section[\texorpdfstring{\gu{દાક્તર} (dāktar)}{દાક્તર (dāktar)}]{\gu{દાક્તર} (dāktar) — doctor}

\label{lesson:GU-C224-daktar}



\begin{quote}
Before the new one: say the Gujarati for common cold, then the Gujarati for to sneeze.
\end{quote}

\subsection*{You'll want to know: \gu{દાક્તર}}
\textbf{\gu{દાક્તર}} — \emph{dāktar} — \textquotedblleft{}doctor\textquotedblright{}.

\begin{quote}
\textbf{\gu{શું} \gu{દાક્તર} \gu{આજે} \gu{અહીં} \gu{છે}?} — \emph{śũ dāktar āje ahī̃ che?} — \textquotedblleft{}Is the doctor here today?\textquotedblright{}
\end{quote}

A masculine noun.

\begin{cousinweb}
\textbf{\gu{દાક્તર}} (\emph{dāktar}) is a loan, usually traced to English \emph{doctor}; EtymDB lists Persian \emph{dāktar} as the go-between.
\end{cousinweb}

\begin{grammarlens}[title={What you've built}]
One more word for everyday life at A2.
\end{grammarlens}

\subsection*{Your turn}
\begin{itemize}
\raggedright
  \item \emph{Say it:} \emph{dāktar}
  \item \emph{Say it:} \emph{dāktar}, once more
  \item \emph{Say it:} \emph{śũ dāktar āje ahī̃ che?} — Is the doctor here today?
  \item \emph{Recall:} say the Gujarati for common cold, then the Gujarati for to sneeze, then say \emph{dāktar} again
\end{itemize}

\begin{tcolorbox}[breakable,colback=teal!4,colframe=teal!35!black,title={Before you move on}]
What does \gu{દાક્તર} mean? (\textquotedblleft{}Doctor\textquotedblright{}.) Say it once more.
\end{tcolorbox}
\section[\texorpdfstring{\gu{દર્દી} (dardī)}{દર્દી (dardī)}]{\gu{દર્દી} (dardī) — patient}

\label{lesson:GU-C224-dardi}



\begin{quote}
Before the new one: say the Gujarati for to sneeze, then the Gujarati for doctor.
\end{quote}

\subsection*{You'll want to know: \gu{દર્દી}}
\textbf{\gu{દર્દી}} — \emph{dardī} — \textquotedblleft{}patient\textquotedblright{}.

\begin{quote}
\textbf{\gu{શું} \gu{દર્દી} \gu{ઘરે} \gu{છે}?} — \emph{śũ dardī ghare che?} — \textquotedblleft{}Is the patient at home?\textquotedblright{}
\end{quote}

A man or a woman.

\begin{cousinweb}
\textbf{\gu{દર્દી}} (\emph{dardī}) is built on \emph{dard}, \textquotedblleft{}pain\textquotedblright{}, a loan usually traced to Persian \emph{dard}, with the suffix \emph{-ī}, \textquotedblleft{}one who has\textquotedblright{}. CDIAL has no entry for it.
\end{cousinweb}

\begin{grammarlens}[title={What you've built}]
One more word for everyday life at A2.
\end{grammarlens}

\subsection*{Your turn}
\begin{itemize}
\raggedright
  \item \emph{Say it:} \emph{dardī}
  \item \emph{Say it:} \emph{dardī}, once more
  \item \emph{Say it:} \emph{śũ dardī ghare che?} — Is the patient at home?
  \item \emph{Recall:} say the Gujarati for to sneeze, then the Gujarati for doctor, then say \emph{dardī} again
\end{itemize}

\begin{tcolorbox}[breakable,colback=teal!4,colframe=teal!35!black,title={Before you move on}]
What does \gu{દર્દી} mean? (\textquotedblleft{}Patient\textquotedblright{}.) Say it once more.
\end{tcolorbox}
\section[\texorpdfstring{\gu{નાડી} (nāḍī)}{નાડી (nāḍī)}]{\gu{નાડી} (nāḍī) — pulse}

\label{lesson:GU-C224-nadi}



\begin{quote}
Before the new one: say the Gujarati for doctor, then the Gujarati for patient.
\end{quote}

\subsection*{You'll want to know: \gu{નાડી}}
\textbf{\gu{નાડી}} — \emph{nāḍī} — \textquotedblleft{}pulse\textquotedblright{}.

\begin{quote}
\textbf{\gu{શું} \gu{દાક્તર} \gu{નાડી} \gu{જુએ} \gu{છે}?} — \emph{śũ dāktar nāḍī jue che?} — \textquotedblleft{}Is the doctor taking the pulse?\textquotedblright{}
\end{quote}

A feminine noun.

\begin{cousinweb}
\textbf{\gu{નાડી}} (\emph{nāḍī}) comes from Sanskrit \emph{nāḍī}, \textquotedblleft{}tubular stalk, tubular organ, tube\textquotedblright{}. CDIAL lists Gujarati \emph{nāṛī}, \textquotedblleft{}pulse\textquotedblright{}, and \emph{nāṛ}, \textquotedblleft{}vein, pulse\textquotedblright{} (CDIAL 7047).
\end{cousinweb}

\begin{grammarlens}[title={What you've built}]
One more word for everyday life at A2.
\end{grammarlens}

\subsection*{Your turn}
\begin{itemize}
\raggedright
  \item \emph{Say it:} \emph{nāḍī}
  \item \emph{Say it:} \emph{nāḍī}, once more
  \item \emph{Say it:} \emph{śũ dāktar nāḍī jue che?} — Is the doctor taking the pulse?
  \item \emph{Recall:} say the Gujarati for doctor, then the Gujarati for patient, then say \emph{nāḍī} again
\end{itemize}

\begin{tcolorbox}[breakable,colback=teal!4,colframe=teal!35!black,title={Before you move on}]
What does \gu{નાડી} mean? (\textquotedblleft{}Pulse\textquotedblright{}.) Say it once more.
\end{tcolorbox}
\section[\texorpdfstring{\gu{ધબકવું} (dhabakvũ)}{ધબકવું (dhabakvũ)}]{\gu{ધબકવું} (dhabakvũ) — to throb, to beat}

\label{lesson:GU-C224-dhabakvu}



\begin{quote}
Before the new one: say the Gujarati for patient, then the Gujarati for pulse.
\end{quote}

\subsection*{You'll want to know: \gu{ધબકવું}}
\textbf{\gu{ધબકવું}} — \emph{dhabakvũ} — \textquotedblleft{}to throb, to beat\textquotedblright{}.

\begin{quote}
\textbf{\gu{શું} \gu{તેનું} \gu{હૃદય} \gu{ધબકે} \gu{છે}?} — \emph{śũ tenũ hṛday dhabke che?} — \textquotedblleft{}Is his heart beating?\textquotedblright{}
\end{quote}

\begin{cousinweb}
The origin of \textbf{\gu{ધબકવું}} (\emph{dhabakvũ}) is not settled. It is probably imitative of the beat, like \emph{dhaṛakvũ}, \textquotedblleft{}to throb\textquotedblright{}, which CDIAL lists under \emph{*dhaḍ}, \textquotedblleft{}throb, beat\textquotedblright{} (CDIAL 6711).
\end{cousinweb}

\begin{grammarlens}[title={What you've built}]
One more word for everyday life at A2.
\end{grammarlens}

\subsection*{Your turn}
\begin{itemize}
\raggedright
  \item \emph{Say it:} \emph{dhabakvũ}
  \item \emph{Say it:} \emph{dhabakvũ}, once more
  \item \emph{Say it:} \emph{śũ tenũ hṛday dhabke che?} — Is his heart beating?
  \item \emph{Recall:} say the Gujarati for patient, then the Gujarati for pulse, then say \emph{dhabakvũ} again
\end{itemize}

\begin{tcolorbox}[breakable,colback=teal!4,colframe=teal!35!black,title={Before you move on}]
What does \gu{ધબકવું} mean? (\textquotedblleft{}To throb, to beat\textquotedblright{}.) Say it once more.
\end{tcolorbox}
\section[\texorpdfstring{\gu{શ્વાસ} (śvās)}{શ્વાસ (śvās)}]{\gu{શ્વાસ} (śvās) — breath}

\label{lesson:GU-C224-shvas}



\begin{quote}
Before the new one: say the Gujarati for pulse, then the Gujarati for to throb.
\end{quote}

\subsection*{You'll want to know: \gu{શ્વાસ}}
\textbf{\gu{શ્વાસ}} — \emph{śvās} — \textquotedblleft{}breath\textquotedblright{}.

\begin{quote}
\textbf{\gu{શું} \gu{તે} \gu{શ્વાસ} \gu{લે} \gu{છે}?} — \emph{śũ te śvās le che?} — \textquotedblleft{}Is he breathing?\textquotedblright{}
\end{quote}

A masculine noun. With \textbf{\gu{લેવું}} it means to breathe.

\begin{cousinweb}
\textbf{\gu{શ્વાસ}} (\emph{śvās}) is Sanskrit \emph{śvāsa}, \textquotedblleft{}breath, breathing\textquotedblright{}, taken over whole as a learned loan, a tatsama. The inherited form of the same word is Gujarati \emph{sās}, \textquotedblleft{}breath\textquotedblright{} (CDIAL 12769).
\end{cousinweb}

\begin{grammarlens}[title={What you've built}]
One more word for everyday life at A2.
\end{grammarlens}

\subsection*{Your turn}
\begin{itemize}
\raggedright
  \item \emph{Say it:} \emph{śvās}
  \item \emph{Say it:} \emph{śvās}, once more
  \item \emph{Say it:} \emph{śũ te śvās le che?} — Is he breathing?
  \item \emph{Recall:} say the Gujarati for pulse, then the Gujarati for to throb, then say \emph{śvās} again
\end{itemize}

\begin{tcolorbox}[breakable,colback=teal!4,colframe=teal!35!black,title={Before you move on}]
What does \gu{શ્વાસ} mean? (\textquotedblleft{}Breath\textquotedblright{}.) Say it once more.
\end{tcolorbox}
\section[(dialogue)]{Words from chapters 219–224 — a review}

\label{lesson:GU-R224-review-219-224}



\begin{quote}
Say the Gujarati for to thunder and to shine.
\end{quote}

\subsection*{Your turn}
Say these aloud:

\begin{itemize}
\raggedright
  \item to thunder, to shine, to cloud over, hailstones, unseasonal rain
  \item to flow, to overflow, flood, to recede, to clear up
  \item drought, to fall off, fire, to catch fire, earthquake
  \item planet earth, atmosphere, nature, high tide, low tide
  \item health, body, to ache, common cold, to sneeze
  \item doctor, patient, pulse, to throb, breath
\end{itemize}

\begin{tcolorbox}[breakable,colback=teal!4,colframe=teal!35!black,title={Before you move on}]
To thunder? (\textbf{\gu{ગાજવું}}.) To shine? (\textbf{\gu{ચમકવું}}.) To cloud over? (\textbf{\gu{છવાવું}}.) Hailstones? (\textbf{\gu{કરા}}.) Unseasonal rain? (\textbf{\gu{માવઠું}}.) To flow? (\textbf{\gu{વહેવું}}.) To overflow? (\textbf{\gu{છલકાવું}}.) Flood? (\textbf{\gu{પૂર}}.) To recede? (\textbf{\gu{ઓસરવું}}.) To clear up? (\textbf{\gu{ઊઘડવું}}.) Drought? (\textbf{\gu{દુકાળ}}.) To fall off? (\textbf{\gu{ખરવું}}.) Fire? (\textbf{\gu{આગ}}.) To catch fire? (\textbf{\gu{સળગવું}}.) Earthquake? (\textbf{\gu{ભૂકંપ}}.) Planet earth? (\textbf{\gu{પૃથ્વી}}.) Atmosphere? (\textbf{\gu{વાતાવરણ}}.) Nature? (\textbf{\gu{કુદરત}}.) High tide? (\textbf{\gu{ભરતી}}.) Low tide? (\textbf{\gu{ઓટ}}.) Health? (\textbf{\gu{તબિયત}}.) Body? (\textbf{\gu{શરીર}}.) To ache? (\textbf{\gu{દુખવું}}.) Common cold? (\textbf{\gu{શરદી}}.) To sneeze? (\textbf{\gu{છીંકવું}}.) Doctor? (\textbf{\gu{દાક્તર}}.) Patient? (\textbf{\gu{દર્દી}}.) Pulse? (\textbf{\gu{નાડી}}.) To throb? (\textbf{\gu{ધબકવું}}.) Breath? (\textbf{\gu{શ્વાસ}}.)
\end{tcolorbox}
